% supplementary.tex

\section{Supplementary methods}

\subsection{Additional details for the diverse MD panel}

The diverse panel was assembled from an archived SAbDab summary and the
corresponding IMGT-renumbered structures \citep{Dunbar2014SAbDab}. For each PDB
identifier, the workflow took the first heavy-chain entry whose chain identifier
was present in the structure file. After availability checks, this deterministic
selection yielded 1143 structures with a usable sequence, topology, trajectory,
and native-contact label.

The selected heavy chain was kept as an uncapped protein at pH 7.4, solvated in
OPC water with 15~\AA{} padding and 0.15~M NaCl, and described with ff19SB
\citep{Tian2020ff19SB,Izadi2014OPC}. The 300-K branch comprised 1~ns NVT and
1~ns NPT with a 2-fs timestep. Two 1-ns restrained NVT stages at 400~K preceded
100~ns of unrestrained 400-K NVT production. Production used hydrogen-mass
repartitioning to 4~amu, a 4-fs timestep, and output every 100~ps. OpenMM used
PME electrostatics with a 1.0-nm cutoff, constraints on bonds to hydrogen, and a
Langevin middle integrator with 1~ps$^{-1}$ friction
\citep{Eastman2017OpenMM}.

The diverse-panel contact calculation used backbone heavy-atom pairs separated
by more than three residues, within 0.45~nm in the reference structure, and with
at least one atom in Asp, Glu, Gln, Asn, Arg, Lys, or His. For trajectories with
more than 300 frames, the label was averaged over the larger of the final 300
frames and the final one third of the trajectory; all frames were used for
shorter trajectories. Supplementary Fig.~S1 shows the relation between sequence
length and the resulting raw $Q$ value.

\subsection{FEP provenance and charge-change check}

The retained FEP tables contain Ala, Asp, Ile, and Gln scans. Row-level
provenance was checked against the archived scan outputs; the available Phe
scans were not included.

For charge-changing mutations in a cubic 70-\AA{} box and solvent dielectric
78, we tested an analytical periodic net-charge correction of
$0.086(\Delta q)^2$ kcal mol$^{-1}$. We applied the correction before repeating
the sign reversal, robust scaling, Yeo--Johnson transform, and min--max scaling
used for the model-facing labels (Supplementary Fig.~S2b).

\renewcommand{\thefigure}{S\arabic{figure}}

\begin{figure}[p]
\centering
\includegraphics[width=\textwidth]{supp_fig01_data_and_design}
\caption{Sequence length and native-contact $Q$ in the heterogeneous MD panel.
Hexagons show row density, and dashed lines mark the fixed sequence lengths of
the 1MEL and 4IDL matched mutation scans. The observed association is
descriptive and does not establish that length caused the transfer difference.}
\label{fig:supp_data_design}
\end{figure}

\begin{figure}[p]
\centering
\includegraphics[width=\textwidth]{supp_fig02_transfer_controls}
\caption{Robustness controls for the FEP result.
\textbf{(a)} Paired FEP-minus-Tm-only MAE for fine-tuned ESM2 encoders. Negative
values mean lower Tm error with FEP. The 8M point is the staged-selected
five-seed model; the 35M and 650M points are exploratory fixed-configuration
ten-seed controls and were not subjected to the same search. Whiskers are
nominal 90\% paired bootstrap intervals over the 396 test proteins.
\textbf{(b)} Change in each scaled FEP label after the analytical periodic
net-charge correction at dielectric 78.}
\label{fig:supp_transfer_controls}
\end{figure}
